% 1-WL 在三点路径 1-2-3 上：c^(0)=(A,A,A)；签名 (A,{{A}})->B，(A,{{A,A}})->C；c^(1)=(B,C,B)；
% 签名 (B,{{C}})->D，(C,{{B,B}})->E；c^(2)=(D,E,D)，分区 {{1,3},{2}} 不再细化。
\begin{tikzpicture}[font=\figurefont\fontsize{8}{10}\selectfont,>=Stealth,
  vtx/.style={circle,draw=FigureStroke,line width=.8pt,minimum size=6.2mm,inner sep=0pt,text=FigureInk},
  edge/.style={draw=FigureStroke,line width=.85pt},
  idx/.style={font=\FigureNoteFont\fontsize{7.5}{9}\selectfont,text=FigureMuted},
  note/.style={font=\FigureNoteFont\fontsize{7.5}{9}\selectfont,text=FigureMuted}]
  \foreach \r/\t/\la/\lb/\fa/\fb in {0/0/A/A/FigureBlueFill/FigureBlueFill,1/1/B/C/FigureGreenFill/FigureYellowFill,2/2/D/E/FigureGreenFill/FigureYellowFill}{
    \pgfmathsetmacro{\y}{-1.15*\r}
    \node[anchor=west] at (-1.3,\y) {$c^{(\t)}$};
    \node[vtx,fill=\fa] (p\r1) at (0,\y) {$\la$};
    \node[vtx,fill=\fb] (p\r2) at (1.0,\y) {$\lb$};
    \node[vtx,fill=\fa] (p\r3) at (2.0,\y) {$\la$};
    \draw[edge] (p\r1)--(p\r2)--(p\r3);
  }
  \foreach \i in {1,2,3} \node[idx,above=0pt of p0\i] {顶点 \i};
  % 签名
  \node[anchor=west] at (2.8,-0.35) {端点 $(A,\{\!\{A\}\!\})\mapsto B$，中点 $(A,\{\!\{A,A\}\!\})\mapsto C$};
  \node[anchor=west] at (2.8,-1.5) {端点 $(B,\{\!\{C\}\!\})\mapsto D$，中点 $(C,\{\!\{B,B\}\!\})\mapsto E$};
  \node[anchor=west,note] at (2.8,-2.3) {名称变了，分区仍是 $\{1,3\}$、$\{2\}$：停止};
  \draw[->,FigureMuted,line width=.6pt] (2.6,-0.15)--(2.6,-0.85);
  \draw[->,FigureMuted,line width=.6pt] (2.6,-1.3)--(2.6,-2.0);
\end{tikzpicture}
